\section{Flux as deployed today}
\label{sec:deployed-today}\label{sec:baseline}

Most descriptions of \Flux in circulation — including earlier drafts of this paper's own
introduction — are anchored to whatever fleet size was true when someone last looked. The network
does not stand still, and a reader who has not looked recently will underestimate both its scale
and its heterogeneity: how many nodes, how much collateral is actually locked behind them, how many
applications are running, and how old some of those applications' specifications are. This section
replaces impression with a single dated measurement, taken directly from the public API surface
that any third party can query, and states the numbers that the rest of this paper treats as the
baseline against which every later claim — about consensus, about pricing, about the target
architecture — is measured.

\subsection{Methodology}
\label{sec:deployed-today-methodology}

Two conventions matter for reading every figure in this paper. First, the snapshot was taken over a
short window on 2026-09-03 rather than at a single instant: the API captures correspond to chain
height 2{,}917{,}115, while the supply figures computed by the emission model are stated at height
2{,}912{,}015, roughly 42 hours earlier at the \SI{30}{\second} target spacing. \textbf{Every figure
therefore names the height it was taken at}, and heights, not wall-clock times, are the primary
index. Second, the \PoN transition quadrupled the block rate, so block-denominated durations carry a
fork-aware $4\times$ multiplier: a constant of 22{,}000 blocks corresponds to an effective
88{,}000 blocks after activation, and this paper gives the constant together with the multiplier
rather than silently reporting one of the two.


Every number in this section is computed from a live query against four endpoints of the public
\texttt{api.\allowbreak{}runonflux.\allowbreak{}io} API, retrieved on 2026-09-03:
\texttt{/\allowbreak{}daemon/\allowbreak{}getzelnodecount} (fleet totals by tier and address family),
\texttt{/\allowbreak{}daemon/\allowbreak{}listzelnodes} (the full per-node record set, from which tier collateral, node
tenure, and the tier table below are recomputed directly from the records rather than trusted from
the summary counts), \texttt{/\allowbreak{}apps/\allowbreak{}globalappsspecifications} (the full set of currently registered
application specifications), and \texttt{/\allowbreak{}apps/\allowbreak{}deploymentinformation} (the height-indexed pricing
table and deployment limits currently enforced by \FluxOS). No other data source is used in this
section. The snapshot is reproducible: \texttt{scripts/\allowbreak{}05-network-analysis.\allowbreak{}py} issues the same four
queries and recomputes every table and figure below from the raw response bodies, so a reader can
re-run it against the live network and obtain a different, equally valid snapshot for a later date.
Every figure, table, and number below carries its endpoint and the 2026-09-03 retrieval date
explicitly, following the citation convention used throughout this paper
(\shipped; \srcmeasured{api.runonflux.io}{2026-09-03}). This methodology is stated in full here
because later evaluation (\S\ref{sec:evaluation}) reuses it without repeating it.

\subsection{The fleet}
\label{sec:deployed-today-fleet}

At retrieval time the network comprised \textbf{6,489} \FluxNode{}s, all reported as
\texttt{stable} \shipped\ \srcmeasured{api.runonflux.io/\allowbreak{}daemon/\allowbreak{}getzelnodecount}{2026-09-03}. Of
these, \textbf{2,322} report an IPv4 endpoint; \textbf{0} report IPv6, and \textbf{0} report an
onion (Tor) endpoint \shipped\ \srcmeasured{api.runonflux.io/\allowbreak{}daemon/\allowbreak{}getzelnodecount}{2026-09-03}.
Recomputing tier membership and collateral directly from the \texttt{/\allowbreak{}daemon/\allowbreak{}listzelnodes} record
set (rather than trusting the summary counter) gives the following:

\begin{table}[H]
\centering
\begin{tabular}{lrrr}
\hline
Tier & Nodes & Collateral per node (\flux) & Tier collateral (\flux) \\
\hline
Cumulus & 3,247 & 1,000 & 3,247,000 \\
Nimbus  & 1,615 & 12,500 & 20,187,500 \\
Stratus & 1,627 & 40,000 & 65,080,000 \\
\hline
\textbf{Total} & \textbf{6,489} & --- & \textbf{88,514,500} \\
\hline
\end{tabular}
\caption{Fleet composition and locked collateral by tier, recomputed from per-node records
\shipped\ \srcmeasured{api.runonflux.io/\allowbreak{}daemon/\allowbreak{}listzelnodes}{2026-09-03}.}
\label{tab:fleet-tiers}
\end{table}


\paragraph{What that fleet amounts to, and what it is.} Two aggregates matter and they are not the
same number. \FluxOS{} declares each tier's allocation in \texttt{flux\allowbreak{}Specifics}, and
records against each line, in the configuration's own comments, how much of it is available to
applications after node overhead
\shipped~\srccode{flux/\allowbreak{}ZelBack/\allowbreak{}config/\allowbreak{}default.js}{380--395}. Multiplying both by the
census gives

\begin{center}
\small
\begin{tabular}{@{}lrrr@{}}
\toprule
& \textbf{vCPU} & \textbf{RAM} & \textbf{Storage} \\
\midrule
Total allocation & \num{51940} & \SI{166.4}{\tebi\byte} (\num{170426}\,GB) & \SI{2.86}{\peta\byte} (\num{2856700}\,GB) \\
Usable by applications & \num{45451} & \SI{153.8}{\tebi\byte} (\num{157448}\,GB) & \SI{2.60}{\peta\byte} (\num{2597140}\,GB) \\
\bottomrule
\end{tabular}
\end{center}

\noindent\shipped \srcdoc{scripts/23-network-capacity.py}. The CPU figures count \emph{threads}, not
physical cores: \texttt{flux\allowbreak{}Specifics} states Cumulus as 40 tenths of a vCPU, and
\fluxbench{} states the same machine as 2 cores and 4 threads
\srccode{fluxbench/\allowbreak{}src/\allowbreak{}fluxnode/\allowbreak{}benchmarks.h}{58--84} --- one machine counted two
ways, and the smaller number is not the one \FluxOS{} allocates against.

Both rows are floors rather than measurements. The tier line is what a node must provide to be
confirmed, not a declaration of what it carries, and nothing in the protocol caps a node above its
tier; the true capacity is larger by an amount this survey does not measure. Storage carries a third
and stricter figure: the runtime admission check computes usable disk as
$0.95 \times \text{total} - 60 - 20$ GB
\shipped~\srccode{flux/\allowbreak{}ZelBack/\allowbreak{}src/\allowbreak{}services/\allowbreak{}appRequirements/\allowbreak{}hwRequirements.js}{385},
which for a \num{220}\,GB Cumulus admits \num{129}\,GB rather than the \num{180}\,GB its tier line
records, and across the fleet gives \SI{2.19}{\peta\byte} --- \SI{77}{\percent} of the total
allocation. \decided{both figures are published, because they measure different things and neither substitutes
for the other \srcintent{design intake}. The tier line states what the fleet was
specified to provide; the runtime check states what a deployment can actually obtain. A reader sizing
a workload wants the second; a reader assessing what the network was built to hold wants the first.
The \SI{23}{\percent} gap between them is a real property of the deployed configuration, reported
rather than resolved away.}

On that basis \Flux{} is, to the author's knowledge, the \textbf{largest decentralized compute
network by independently-operated nodes} \srcintent{design intake}. The paper
states the measured basis rather than the bare superlative, and does not attempt a comparison table
against other networks: no such measurement was made for this paper, and a claim of that shape is
only worth as much as the census behind it. What is measured is the census above --- \num{6489}
nodes, \num{88514500}\flux{} of operator-posted collateral, every node independently operated and
independently collateralised.

\notestablished{the exact all-time-high fleet size, though it is narrower than an open question.
Secondary reports put the peak at over \num{14000} nodes, with about \num{13500} in September 2024
and over \num{10000} in mid-2025 against the \num{6489} measured here
\srcintent{design intake}; these are marketing and review pages rather than primary
sources and none is authoritative on its own. One independent check makes the \num{13500} figure
credible: the same reports state more than \num{107000} CPU cores at that node count, and scaling
this section's own per-tier method to \num{13500} nodes gives \num{108058} vCPU --- a
\SI{1}{\percent} agreement, which corroborates both the reported size and the threads-not-cores
basis used above. The exact number is recoverable without any statistics service, since \FluxNode{}
start and confirmation transactions are on-chain and the confirmed set at any past height can be
reconstructed by replaying them; this paper states that method rather than running it.

The contraction is the more consequential fact. \textbf{The fleet stands at roughly half its peak},
which is why the scale claim above is made against the measured present and not against a historical
high --- and which bears directly on \Cref{sec:pnr}'s open question about exit composition under a
demand shortfall: the contraction that section reasons about in the abstract has already happened at
scale, and whether the operators who left were disproportionately single-node holders is still
unmeasured.}


The total of \textbf{88,514,500}\flux is the sum of collateral each operator has locked to run a
node at that tier; it is not a measure of trading volume, market value, or any other quantity this
paper's register excludes. Stratus nodes are the smallest population by count (1,627 of 6,489, or
25.1\%) but hold the largest share of locked collateral (65,080,000 of 88,514,500\flux, or 73.5\%)
\shipped\ \srcmeasured{api.runonflux.io/\allowbreak{}daemon/\allowbreak{}listzelnodes}{2026-09-03}.

\subsection{Node tenure}
\label{sec:deployed-today-tenure}

Measuring blocks elapsed since each node's \texttt{added\_height} against the chain tip at
retrieval time gives a median tenure of \textbf{193,620} blocks (approximately 67 days at the
current 30-second target spacing), a 10th-percentile tenure of \textbf{12,909} blocks, a
90th-percentile tenure of \textbf{881,770} blocks, and a maximum of \textbf{1,648,700} blocks
\shipped\ \srcmeasured{api.runonflux.io/\allowbreak{}daemon/\allowbreak{}listzelnodes}{2026-09-03}. The distance between the
10th and 90th percentiles spans roughly two orders of magnitude, which is consistent with a fleet
that is simultaneously accumulating a long tail of multi-year operators and continuing to admit new
entrants rather than having converged to a stable, unchanging membership \srcinferred.

\subsection{Deployed applications}
\label{sec:deployed-today-apps}

The global application registry held \textbf{1,195} application specifications at retrieval time,
comprising \textbf{1,371} components (compose entries) across \textbf{7,486} total requested
instances \shipped\ \srcmeasured{api.runonflux.io/\allowbreak{}apps/\allowbreak{}globalappsspecifications}{2026-09-03}.
Summed across all requested instances, the committed resource footprint specified by these
applications is \textbf{8,972.0} vCPU-equivalents, \textbf{17,544,229}~MB of RAM (17,133.0~GiB),
and \textbf{160,107}~GB of disk
\shipped\ \srcmeasured{api.runonflux.io/\allowbreak{}apps/\allowbreak{}globalappsspecifications}{2026-09-03}. Instances
requested per application have a median of \textbf{2}, a mean of \textbf{6.26}, a maximum of
\textbf{100}, and a minimum of \textbf{0}; \textbf{439} applications (36.7\%) sit exactly at the
network-wide instance floor of 3
\shipped\ \srcmeasured{api.runonflux.io/\allowbreak{}apps/\allowbreak{}globalappsspecifications}{2026-09-03}.

These figures are \emph{requested} resources as declared in each application's specification, not
measured utilisation. No public endpoint reports realised CPU, memory, disk, or bandwidth
utilisation on running containers across the fleet; this paper does not estimate one, and the
absence is recorded as an open problem rather than papered over with an inferred figure
(\S\ref{sec:limitations}).

\subsection{Specification version distribution}
\label{sec:deployed-today-specversions}

The 1,195 deployed applications are governed by seven distinct \appspec versions, all live
simultaneously at retrieval time:

\begin{table}[H]
\centering
\begin{tabular}{lrr}
\hline
Version & Applications & Share \\
\hline
\specv{2} & 2   & 0.2\% \\
\specv{3} & 27  & 2.3\% \\
\specv{4} & 11  & 0.9\% \\
\specv{5} & 3   & 0.3\% \\
\specv{6} & 70  & 5.9\% \\
\specv{7} & 209 & 17.5\% \\
\specv{8} & 873 & 73.1\% \\
\hline
\textbf{Total} & \textbf{1,195} & \textbf{100\%} \\
\hline
\end{tabular}
\caption{Distribution of deployed applications by \appspec version
\shipped\ \srcmeasured{api.runonflux.io/\allowbreak{}apps/\allowbreak{}globalappsspecifications}{2026-09-03}.}
\label{tab:specversions}
\end{table}

Seven \appspec versions are live simultaneously, spanning from \specv{2} (registered years before
\specv{8} became the newest enforced version) to \specv{8} itself. No version among them has ever
been forced out of the running population by the introduction of a later one \srcinferred — this
reading follows from the observed coexistence rather than from a direct record of deprecation
history, and it is stated hedged for that reason. \S10 examines the version-gating mechanism that
produces this coexistence, and the pending \specv{9} redesign that would eventually extend it to
eight.

\subsection{Pricing as deployed}
\label{sec:deployed-today-pricing}

The rates \FluxOS charges for deployment resources are set by a height-indexed table served from
\texttt{/\allowbreak{}apps/\allowbreak{}deploymentinformation}. At retrieval time the table held six recorded price states:
an initial schedule in force from height $-1$ (i.e., from genesis of the pricing mechanism) and five
subsequent revisions.

\begin{table}[H]
\centering
\begin{tabular}{rrrrrrrr}
\hline
Height & cpu & ram & hdd & minPrice & port & scope & staticip \\
\hline
$-1$       & 3.00 & 1.00 & 0.500 & 1.00 & 2.0 & 6.0 & 3.0 \\
983,000    & 0.30 & 0.10 & 0.050 & 0.10 & 2.0 & 6.0 & 3.0 \\
1,004,000  & 0.06 & 0.02 & 0.010 & 0.01 & 2.0 & 6.0 & 3.0 \\
1,288,000  & 0.15 & 0.05 & 0.020 & 0.01 & 2.0 & 6.0 & 3.0 \\
1,594,832  & 0.15 & 0.05 & 0.020 & 0.01 & 1.5 & 6.0 & 3.0 \\
1,597,156  & 0.03 & 0.01 & 0.004 & 0.01 & 0.4 & 0.8 & 0.4 \\
\hline
\end{tabular}
\caption{Height-scheduled deployment pricing table, in \flux per priced unit
\shipped\ \srcmeasured{api.runonflux.io/\allowbreak{}apps/\allowbreak{}deploymentinformation}{2026-09-03}.}
\label{tab:pricing-schedule}
\end{table}

The rates in force at retrieval time are the row at height 1,597,156. Across the six recorded rows,
the \texttt{cpu} rate has fallen from 3 to 0.03\flux per unit; this is stated as a fact about the
schedule as recorded, with no comment on cost, value, or trend implied.
Payments for deployment and renewal are sent to \texttt{t3Nryf\allowbreak{}AQLGe\allowbreak{}Fs9j\allowbreak{}Eoeqsxm\allowbreak{}BN2QLRa\allowbreak{}RKFLUX}
\shipped\ \srcmeasured{api.runonflux.io/\allowbreak{}apps/\allowbreak{}deploymentinformation}{2026-09-03}. The same endpoint
reports \texttt{minimum\allowbreak{}Instances} of \textbf{3}, \texttt{maximum\allowbreak{}Instances} of \textbf{100},
\texttt{blocksLasting} of \textbf{22,000} blocks, \texttt{max\allowbreak{}Blocks\allowbreak{}Allowance} of
\textbf{1,056,000} blocks, and \texttt{maxImageSize} of \textbf{5,000,000,000} bytes (5~GB)
\shipped\ \srcmeasured{api.runonflux.io/\allowbreak{}apps/\allowbreak{}deploymentinformation}{2026-09-03}.

\subsection{The multi-chain footprint}
\label{sec:deployed-today-parallel}

Beyond the main chain, \Flux is represented today by a parallel asset on ten separate chains, each
carrying its own contract address, mint, or module identifier:

\begin{table}[H]
\centering
\begin{tabular}{llp{6cm}}
\hline
Chain & Identifier & Provenance \\
\hline
Ethereum & \texttt{0x720cd16b011b98\allowbreak{}7da3518fbf38c307\allowbreak{}1d4f0d1495} (ERC-20) & \srccode{holder-snapshot/\allowbreak{}index.js}{22} \\
BNB Chain & \texttt{0xaff9084f237458\allowbreak{}5879e8b434c399e2\allowbreak{}9e80cce635} (BEP-20) & \srccode{holder-snapshot/\allowbreak{}index.js}{23} \\
Polygon & \texttt{0xA2bb7A68c46b53\allowbreak{}f6BbF6cC91C865Ae\allowbreak{}247A82E99B} & \srccode{holder-snapshot/\allowbreak{}index.js}{26} \\
Avalanche C-Chain & \texttt{0xc4B06F17ECcB22\allowbreak{}15a5DBf042C67210\allowbreak{}1Fc20daF55} & \srccode{holder-snapshot/\allowbreak{}index.js}{24} \\
Base & \texttt{0xb008bdcf9cdff9\allowbreak{}da684a190941dc3d\allowbreak{}ca8c2cdd44} & \srccode{flux-supply-tracker/\allowbreak{}index.html}{346} \\
Solana & \texttt{FLUX1wa2Gmbt\allowbreak{}SB6ZGi2p\allowbreak{}TNb\allowbreak{}VCw3z\allowbreak{}Ee\allowbreak{}Kna\allowbreak{}PCk\allowbreak{}Pt\allowbreak{}FXxq\allowbreak{}Xe} (SPL mint) & \srccode{holder-snapshot/\allowbreak{}index.js}{20} \\
Algorand & \texttt{1029804829} (ASA ID) & \srccode{holder-snapshot/\allowbreak{}index.js}{27} \\
Ergo & \texttt{e8b20745ee9d1881\allowbreak{}7305f32eb2101583\allowbreak{}1a48f02d40980de6\allowbreak{}e849f886dca7f807\allowbreak{}} (token ID) & \srccode{holder-snapshot/\allowbreak{}index.js}{25} \\
Tron & \texttt{TWr6yzuk\allowbreak{}Rw\allowbreak{}Z53HDe3bzc\allowbreak{}C8RCTbi\allowbreak{}Ka4Zzb6} (TRC-20) & \srccode{holder-snapshot/\allowbreak{}index.js}{21} \\
Kadena & \texttt{runonflux.flux} (Pact module) & \srccode{flux-supply-tracker/\allowbreak{}index.html}{391} \\
\hline
\end{tabular}
\caption{Chains carrying a \Flux parallel asset, with identifiers as recorded in the team's own
asset-inventory tooling \shipped.}
\label{tab:pa-overview}
\end{table}

These identifiers are as recorded in the codebases that track them (\texttt{holder-\allowbreak{}snapshot},
\texttt{flux-\allowbreak{}supply-\allowbreak{}tracker}); this survey did not independently re-verify each address against
its chain's own explorer \srcinferred. \S6 takes this footprint as the starting inventory for an
argument that consolidating it reduces audit and reserve-accounting surface.

\subsection{The operational layer}
\label{sec:deployed-today-operational}

Around \FluxOS on every node sits a cluster of smaller, single-purpose daemons that this paper
returns to in later sections: \texttt{flux-\allowbreak{}telemetryd} ships per-container metrics and logs to a
backend the application owner — not \Flux — chooses and credentials, never touching data \FluxOS
has not vouched for \shipped\ \srccode{flux-telemetryd/\allowbreak{}src/\allowbreak{}identity.rs}{1--31,56--74} (\S9);
\texttt{flux-fstrimd} paces filesystem TRIM operations so that routine storage maintenance does
not stall a node's I/O \shipped\ \srccode{flux-fstrimd/\allowbreak{}src/\allowbreak{}main.rs}{1--6}; \texttt{flux-shutdownd} is
building a coordinated, reason-aware graceful-shutdown and upgrade pipeline in place of an
unceremonious container kill, though at present only its signal-and-cleanup stages run in
production and its drain and pre-stop execution stages are implemented but not yet wired to real
events \inprogress\ \srccode{flux-shutdownd/\allowbreak{}crates/\allowbreak{}shutdownd/\allowbreak{}src/\allowbreak{}pipeline/\allowbreak{}state\_machine.rs}{3--7};
and, operated centrally by the Foundation rather than per-node,
\texttt{appsmonitor} tracks on-chain application expiry and sends renewal notifications \shipped\
\srccode{appsmonitor/\allowbreak{}src/\allowbreak{}services/\allowbreak{}appsSubscriptionRenewalMonitor.js}{112--368} (\S16),
\texttt{fluxstorage} holds oversized environment-variable and command payloads that do not fit an
on-chain specification \shipped\ \srccode{fluxstorage/\allowbreak{}src/\allowbreak{}routes.js}{17--67} (\S23), and
\texttt{fluxos-\allowbreak{}release-\allowbreak{}api} serves the FluxOS/ArcaneOS upgrade bundles that nodes poll for
\shipped\ \srccode{fluxos-release-api/\allowbreak{}src/\allowbreak{}flux\_release\_api/\allowbreak{}api.py}{176--229} (\S23).
\FluxOS's own placement and lifecycle behaviour that these services
support is described in \S9; the centrally-operated members of this list reappear in \S12's
inventory of components the decentralization roadmap must still address.

\begin{figure}[H]
\centering
\begin{tikzpicture}
\begin{axis}[
    ybar,
    width=0.9\linewidth,
    height=5.5cm,
    ylabel={applications},
    xlabel={\appspec version},
    symbolic x coords={2,3,4,5,6,7,8},
    xtick=data,
    nodes near coords,
    nodes near coords align={vertical},
    ymin=0,
]
\addplot table [x index=0, y index=1] {../data/appspec_versions.dat};
\end{axis}
\end{tikzpicture}
\caption{Number of deployed applications by \appspec version, out of 1,195 total
\shipped\ \srcmeasured{api.runonflux.io/\allowbreak{}apps/\allowbreak{}globalappsspecifications}{2026-09-03}. The
distribution is heavily concentrated at \specv{8} (873 of 1,195) while six older versions remain
live.}
\label{fig:appspec-versions}
%% verified 2026-09-04: appspec_versions.dat header is "# version count share", and this plot reads
%% x index=0, y index=1 --- version against count, as captioned. (Build note; not paper text.)
\end{figure}

\subsection{The measured baseline}
\label{sec:deployed-today-baseline}

Three numbers from this section recur throughout the rest of this paper as the measured baseline
against which every later architectural claim is checked: \textbf{6,489} nodes, \textbf{1,195}
deployed applications, and \textbf{88,514,500}\flux of locked collateral
\shipped\ \srcmeasured{api.runonflux.io}{2026-09-03}. Neither the consensus mechanism of \S4, nor
the compensation design of \S7, nor the target architecture of Part~IV should be read as
describing a network smaller or more homogeneous than this.

%% Drain this section's float queue before the next begins.
\FloatBarrier
